% Source: https://artofproblemsolving.com/community/c3h1949688_formula_to_solve_an_equation
% Stable identity: aops-topic-1949688-q01
Add $65$ to both sides, then divide by $2$:
\[
2w^2-65=63\ \Longrightarrow\ 2w^2=128\ \Longrightarrow\ w^2=64 .
\]
Now $w^2=64$ means $w^2-64=0$, that is, $(w-8)(w+8)=0$, so $w=8$ or $w=-8$. The solutions are $\boxed{w=\pm8}$. (Check: $2\cdot64-65=63$.)

In general, if $c>0$ then $x^2=c$ is equivalent to
\[
x^2-c=0\quad\Longleftrightarrow\quad (x-\sqrt c)(x+\sqrt c)=0 ,
\]
using the difference of squares and $(\sqrt c)^2=c$. A product of two real numbers is zero exactly when one of the factors is zero, so $x=\sqrt c$ or $x=-\sqrt c$, and both of these do satisfy the equation. This is not a separate formula to memorise: it is the factorisation of a difference of two squares. (The symbol $\sqrt c$ always means the nonnegative root, which is why the solution $-\sqrt c$ has to be written separately.)
